\section{阅读顺序及算例组织}
\label{sec:01_05}

本书采用由问题描述进入一般模型，再通过数值实例理解约束、通过程序核对实现的阅读顺序。这一组织方式借鉴VPP第三章将模型、解释性算例与求解代码相互对应的论证安排\cite{VPP2020}。阅读时，应先弄清决策对象及其允许的动作，再检查参数代入后的限制，最后解释求解输出；脱离变量含义直接比较目标值，容易把不同问题的结果混为一谈。

第2章从一般集合打包与带容量的候选选择出发，说明资源关联如何限制任务组合，并以十二候选任务算例完整展开目标、变量域和各条资源约束。这个小例子的作用是让读者能够沿着具体候选追踪资源占用，逐项复算收益。第3章接续介绍SCIP、Streaming、Greedy与ILS以及原作者完整流程，把算法操作与已有实验结果联系起来。理解方法差异时，应同时阅读停止规则和失败记录，不能只看有效运行的目标值。

工程部分宜按第4至第6章顺读。第4章建立杆塔、观测任务、车辆、无人机和电池等物理对象，解释运动、成像、道路、天气及能量关系；第5章把这些关系接入任务分配、航次、会合、资源周转和状态接续，形成完整一般模型。第6章承担从符号到数值的转换：参数采用什么单位，连续关系怎样形成有限数值模型，以及离散精度和验收条件怎样约束结果的解释。需要复现区域算例的读者，尤其应把一般模型的适用条件与有限转录的近似条件对应起来。

解释结果时，需要分别核对保存的候选、验收记录和求解停止状态，不能用模型规模替代实际覆盖程度。纯组合实验用于讨论资源耦合和算法表现；巡检组件实验用于核查物理关系及其数值实现。阅读两类结果时，应保持问题范围、可行域和完成标准一致。

算例按数学结构组织。每一种不同模型类型均以一套完整数值实例为展开单位，列明集合成员、参数、变量域、目标函数和全部实例化约束，使读者能够核查每个系数及每个索引对应的对象。若仅改变规模、拓扑或参数，而变量含义和约束结构不变，则按同型条件组织完整参数、结果及约束导出入口，供读者按需取得全部约束。若新增状态、改变决策时序或引入新的约束关系，则需重新说明相应模型，不能把结构变化隐藏在参数表中。

正文公式、数值实例、解的验收和电子源码构成一条对照阅读路径。先从公式确认限制的动作及连接的状态，再在实例中代入参数，随后用保存的决策复算目标并检查约束，最后对照代码中的数据读取、变量建立、约束生成和输出处理。电子源码保留可运行文件及其依赖信息，配合公式与实现的对应说明。阅读代码时还须区分求解器返回状态与验收结论：程序结束不自动意味着方案有效，方案有效也不自动意味着完整覆盖或全局最优。
